\thispagestyle{empty}
\begin{center}
{\fontsize{9}{11}\selectfont Manuscrito de recopilación para transferencia académica\par}
\vspace{6mm}
{\fontsize{23}{27}\selectfont\bfseries Relaciones estructurales\par entre las constantes físicas\par}
\vspace{5mm}
{\fontsize{13}{17}\selectfont Acción, reciprocidad geométrica\par y memoria de las transformaciones\par}
\vspace{8mm}
{\large Oumar Haidara Fall \textperiodcentered\ Rubén Ramos Balsa\par}
\vspace{2mm}
{\fontsize{8}{10}\selectfont Coautoría sin prelación de contribución\par}
\vspace{2mm}
{\fontsize{8}{10}\selectfont Integración y recopilación documental generadas por ChatGPT - Astra.\par}
\end{center}
\vspace{6mm}
\begin{abstract}
Se establecen relaciones de compatibilidad, reciprocidad y reconstrucción
entre las estructuras de acción, masa y respuesta constitutiva obtenidas
mediante Holografía Modular Triádica y Mecánica Dimensional. El desarrollo
parte de la composición de la estructura aritmética de caminos
\emph{All Possible Paths} (APP), la firma ternaria orientada TRIT y el
\emph{Topological Phase Kernel} (TPK), que actualiza y transporta el estado
enriquecido. Esta composición conserva orientación, acarreo, memoria y
retorno. Las constantes intervienen como evaluaciones de las
estructuras previamente construidas; su composición permite determinar qué
propiedades comparten y qué información recupera una observación conjunta.
En la familia constitutiva normalizada se demuestra que la velocidad de
propagación y el funcional de Barbero--Immirzi determinan los autovalores
etiquetados de los dos canales. Los logaritmos de velocidad e impedancia
constituyen el gradiente de un potencial espectral estrictamente cóncavo.
Su Hessiano caracteriza la sensibilidad de la inversión y proporciona cotas
explícitas de estabilidad. El potencial comparte un momento espectral de
profundidad dos con la construcción orientada cuyo valor extremo es Catalán;
la respuesta constitutiva restituye esa familia de manera única con su
condición de frontera en el origen.

La composición de cuartos de giro asociados a formas elípticas conjugadas
produce un transporte relativo unimodular cuyas dilataciones recíprocas
permiten reconstruir el coeficiente adimensional de acción, conservadas la
coordenada de clausura y la orientación. El levantamiento dodecafásico y la
iteración del transporte caracterizan la memoria del orden de las
operaciones. El análisis cronológico distingue el crecimiento de la
información de la tasa entrópica por posición, y la reentrada de memoria
determina su contribución a la respuesta operatoria. El retorno inscrito y
su composición polar determinan una longitud $L_\alpha$. Conservar
conjuntamente acción y área fija $G=c^3L_\alpha^2/\hbar$ en la realización
radial, también bajo variaciones no conmutativas y reducción compatible
de memoria.

La memoria electrónica se restituye mediante invariantes lineales y
espectrales que conservan el orden del registro; el transporte entre
multisecciones admite un criterio explícito de conservación de masa.
La composición térmica reúne capacidades, multiplicidades y memoria
espectral en una traza exacta, y construye su respuesta termométrica
para una lectura marcada con referencia fija. Esta realización enlaza
energía, información, acción y área, conservando las condiciones de
normalización y de observación de cada magnitud.
La composición transportada de las respuestas y la reconstrucción
Higgs--impedancia enlazan los canales constitutivos con el sector
electrodébil. La sustitución electrónica conserva su escala en Fermi,
la cancela en los cocientes angulares y transporta la longitud inversa
mediante una congruencia operatoria, también sin conmutación.

Finalmente, bajo las
condiciones circular y térmica especificadas, una misma coordenada de masa
determina la frecuencia propia, dos longitudes recíprocas de producto
invariante y la ponderación térmica del ciclo. Los resultados reúnen así
producción, compatibilidad y restitución en una descripción matemática
común, con demostraciones explícitas y con la distinción entre cambio de
representación, transporte relativo y realización física.
\end{abstract}

\clearpage
\section{Introducción}
\enlargethispage{-3\baselineskip}
\label{md:introduccion}
La lectura de una constante como valor escalar representa el término final
de una operación. Cuando se conserva la estructura que produce esa lectura,
se hacen accesibles otras preguntas: qué relaciones comparte con las demás
lecturas, qué transformación conserva su orientación y qué parte de la
historia puede restituirse a partir de observaciones conjuntas. El presente
desarrollo sitúa esas preguntas en la continuidad entre la estructura discreta
HMT y sus realizaciones en Mecánica Dimensional.

El origen común conserva las operaciones de APP, el régimen orientado del
TRIT y la composición efectiva del TPK. El retorno de fase se distingue de
la evolución del estado enriquecido. Las coordenadas regionales y su
clausura dodecafásica se transmiten después a los lectores de acción,
incidencia y masa. Cada transmisión conserva los objetos y las condiciones
que utiliza; una salida puede intervenir en la construcción siguiente sin
perder su procedencia. El núcleo inicial expone esa arquitectura con sus
operaciones, prolongaciones y pruebas. Los antecedentes sectoriales que
siguen materializan las dependencias empleadas por las composiciones del
presente trabajo.

La cuestión central puede formularse en tres términos complementarios:
\emph{producción}, \emph{compatibilidad} y \emph{restitución}. La
producción fija el objeto y su genealogía. La compatibilidad determina qué
relaciones deben satisfacer sus distintas lecturas. La restitución establece
qué datos del objeto pueden recuperarse desde ellas. El recorrido inverso
es una propiedad de las lecturas ya construidas y conserva sus bases,
orientaciones y dominios; no reemplaza su orden generador.

El primer eje del manuscrito reúne la acción con el par angular. El
coeficiente adimensional de la sección de acción participa en el
contraángulo; este determina, con el ángulo, los canales que alimentan la
respuesta constitutiva y el funcional de Barbero--Immirzi. La misma sección
interviene en el cociente de retorno de la lectura electrónica. De esta
forma, sustituir las estructuras en sus ecuaciones revela dependencias
compartidas que una enumeración de constantes separadas deja implícitas.
La distinción entre coeficiente y magnitud dimensional es esencial:
restituir una razón de forma y restituir una acción dimensional requieren
datos de referencia diferentes.

El segundo eje estudia la información conservada en la respuesta. La
velocidad y la impedancia normalizadas retienen aspectos distintos del
reparto entre canales. La composición permite demostrar tanto una inversión
global en la familia declarada como una relación diferencial entre ambas
respuestas. Un potencial espectral organiza esa relación, y su Hessiano
controla la sensibilidad de la inversión. Así se distinguen dos propiedades:
que los datos exactos determinen unívocamente el objeto y que una precisión
finita permita reconstruirlo con un error acotado. La familia orientada cuyo
valor extremo es Catalán se incorpora mediante su operación espectral,
conservando sus argumentos y su profundidad.

La composición transportada de las respuestas proporciona su coordenada
aditiva; las firmas angulares de calibre y de Higgs permiten reconstruir
la misma pareja de canales. Los momentos del defecto cíclico completan
esa lectura con cotas de convergencia y una representación de Mellin.
Más adelante, la sección electrónica se sustituye en Fermi y en la escala
de Higgs: la prueba conserva qué factores permanecen, cuáles se cancelan
y cómo se transforma el inverso del operador de masa sin suponer conmutación.

El tercer eje traslada la discusión al orden de las operaciones. El plano
de cuarto de giro del ciclo dodecafásico y las formas conjugadas de la
elipse permiten construir un conmutador explícito. El balance de los
incrementos de fase se cierra mientras la acción sobre el estado conserva
una dilatación recíproca. Dos programas con las mismas multiplicidades de
operaciones pueden producir respuestas diferentes por su orden. Esa
diferencia admite detectores y recuperadores definidos: la traza registra
ciertos invariantes, mientras las respuestas sobre ejes orientados
distinguen sentidos que la traza identifica.

La lectura física se presenta junto con las condiciones de cada
realización. En la sección gravitatoria circular, la longitud reducida
de Compton y el radio gravitatorio reducido se dilatan recíprocamente y
mantienen un área común. En la realización térmica del ciclo, la misma
coordenada estructural interviene como exponente de un peso de equilibrio.
En el transporte elíptico, la conservación de área simpléctica y la
conservación de una energía cuadrática común se examinan mediante sus
operadores respectivos. La eventual realización experimental del programa
relativo exige especificar conjuntamente sistema, referencia, operaciones
disponibles y controlador. Estas condiciones pertenecen al significado
del resultado y permanecen contiguas a sus pruebas.

La normalización gravitatoria se desarrolla antes de utilizar sus
consecuencias escalares. El retorno de Hadamard, la inscripción de incidencias
y la norma local de $\alpha$ producen una sección de longitud. Su
transporte polar y el mismo carácter de ruta producen dos lectores
recíprocos. La energía conserva a su vez la longitud de acción. Al reunir
$R\Lambda=L_\alpha^2I$ con $H\Lambda=\hbar cI$, la constante gravitatoria
queda fijada como coeficiente común de todos los modos positivos de la
realización. El reloj circular se expresa después desde longitud y
velocidad. Esta precedencia permite explicar la normalización sin utilizar
el valor experimental de $G$ ni una longitud de Planck tabulada.

La composición electrónica y la térmica precisan cómo interviene la
memoria en esas relaciones. En la primera, se reconstruye el registro
dodecafásico a partir de su suma, sus diferencias y su orientación;
el espectro distingue órdenes que comparten multiplicidades. En la
segunda, el espectro geométrico ya construido por el transporte de
memoria aporta la razón $1/10$. Las capacidades cronológicas y los
sectores de ocupación determinan una traza marcada que se compone
con un estado de Gibbs. Se especifican por separado la energía
condicionada, la información marcada y la temperatura resultante;
esta elección de observaciones permite demostrar el transductor
termométrico y seguir su relación con la escala de acción y el área.
Las secciones correspondientes conservan los datos de preparación,
las graduaciones y las hipótesis de variación utilizadas.

La exposición desarrolla primero los antecedentes, conserva después la
narrativa explicativa del recorrido y presenta las composiciones completas.
El estudio cronológico de las vacancias, el episodio regional de reiteración
visible y la reentrada de memoria permiten seguir el origen de las
preguntas sobre compatibilidad. Los desarrollos energéticos y gaussianos
se conservan en los apéndices como realizaciones auxiliares, con sus
hipótesis y pruebas completas.
Las conclusiones reúnen las caracterizaciones estructurales obtenidas.
El suplemento digital preserva las fuentes de trabajo literales, los
programas y sus resultados; el manuscrito contiene las demostraciones
empleadas, y distingue su lectura matemática de los controles finitos que
las acompañan.

Esta organización permite formular con precisión la aportación transversal
del trabajo. Una constante conserva su valor numérico, pero su definición
estructural especifica además la operación que lo produce, el dominio en
que actúa y las transformaciones que conserva. Dos magnitudes que comparten
una dependencia pueden, por ello, restringirse mutuamente: conocer una
respuesta y su orientación permite reconstruir un reparto que una lectura
escalar aislada no distingue. El problema inverso se plantea sobre la imagen
ya generada y conserva sus datos de referencia. Así, generación y
reconstrucción se articulan sin invertir la precedencia de la construcción.

La misma distinción ordena la relación entre memoria e información. El
estado presente, el balance acumulado y el operador que responderá a una
continuación son objetos diferentes. El retorno de una fase puede coexistir
con una transformación persistente del estado; dos sucesiones con idénticos
recuentos pueden diferenciarse por el orden de composición. Las pruebas
desarrollan esa diferencia mediante operadores, invariantes y observaciones
recuperadoras. Su interés físico reside en precisar qué conserva cada
descripción y qué operación permite acceder a la información retenida.

El paso a las magnitudes dimensionales se realiza manteniendo las bases y
normalizaciones correspondientes. La reconstrucción de un coeficiente de
acción, la conservación de un área y la caracterización de una temperatura
pertenecen a resultados relacionados, con tipos distintos. La exposición
reúne esos resultados en sus condiciones de validez y muestra cómo se
componen. Las conclusiones explicitan tanto la caracterización obtenida
para cada objeto como la relación que establece con las demás estructuras.
